\clearpage
\SourceLeaf{97}%

\begin{center}
\large [ON THE NATURE AND DEFINITION OF FUNCTIONS.]
\end{center}

\begin{center}
\textsc{Gauss} to \textsc{Bessel}, 18 December 1811.
\end{center}

\noindent\hspace*{\parindent}$\mathord{.}\,\mathord{.}\,\mathord{.}\,\mathord{.}\,\mathord{.}\,\mathord{.}$ First of all, I would ask anyone who wishes to introduce a new function into analysis to explain whether he intends it to be applied strictly to real quantities (real values of the function's argument), regarding the imaginary values of the argument as a mere excrescence --- or whether he agrees with my principle that in the realm of quantities the imaginary quantities $a+b\sqrt{-1}=a+bi$ must be regarded as having the same rights as the real ones. Practical utility is beside the point here. I regard analysis as an independent science that would lose an extraordinary amount of beauty and completeness by relegating those fictitious quantities, and would be obliged at every turn to impose extremely troublesome restrictions on truths that otherwise hold generally. $\mathord{.}\,\mathord{.}\,\mathord{.}\,\mathord{.}\,\mathord{.}\,\mathord{.}\,\mathord{.}\,\mathord{.}\,\mathord{.}$ What, then, should one understand by $\int \varphi x\mathord{.} dx$ for $x=a+bi$? Clearly, if one wishes to start from clear concepts, one must suppose that $x$ passes by infinitesimal increments (each of the form $a+\beta i$) from the value at which the integral is to be 0 to $x=a+bi$, and then all the $\varphi x\mathord{.} dx$ are \textls{summed}. This fixes the meaning completely. But the passage can be made in infinitely many ways. Just as one can picture the \textls{entire realm of all real quantities} by an infinite straight line, so one can picture the \textls{entire realm of all} quantities, real and
